\section{Responses to README Questions}
\label{sec:readme-questions}

The UsingMatrixClass README ends with seventeen questions.
They are grouped below into four themes: profiling and measurement, allocator
design and workload fit, hardware and operating system effects, and software
engineering practice.
Each question is quoted in italics before its response.

\subsection{Profiling and Measurement}
\label{sec:rq-measurement}

\subsubsection{Profiling Scale}
\emph{What do you notice as the profiling scale increases in each test case?}
The performance gap between the custom allocator and regular \texttt{new}
stayed roughly constant as scale increased.
Peak memory barely changed between the short, medium, and long tests.
This suggests the scale parameter repeats the same working set rather than
enlarging it, which means the pool's memory overhead stayed fixed.
Minor page faults grew almost exactly in proportion to elapsed time.
Overall, the choice of test case mattered more than the scale.

\subsubsection{Fixed Random Seed}
\emph{Why is a fixed random seed important for a fair comparison?}
It is important that all versions face the same seed to ensure fairness when
comparing performance.
This ensures that no one version is doing any more or less work than the
others.
A fixed seed also makes results reproducible.

\subsubsection{Execution Order}
\emph{How could execution order affect measured results?}
Execution order matters because each test inherits the machine state left by
whatever ran before it.
That state includes the operating system's page cache, which already holds the
executable and shared libraries after the first test has loaded them, as well
as the CPU's frequency and thermal state.
A test running second is therefore measured under different conditions than
the same test running first: whichever test runs first pays that cost and
whichever runs second gets it for free.

\subsubsection{Repeated Runs and Median Times}
\emph{Why should several runs and median times be used for profiling?}
Several runs and median times should be used for profiling to avoid the data
being skewed by outliers.
This prevents an anomaly that slowed down one run from affecting the profiling
data too much.

\subsubsection{Measurements Beyond Elapsed Time}
\emph{Which measurements should be collected in addition to elapsed time?}
Because the results are collected across multiple teams, the environment must
be recorded alongside the timing.
This includes CPU model, core count, operating system version, and compiler
version.
Additionally, each run records CPU utilization and peak memory usage.
Each run also records minor page faults, which showed that regular
\texttt{new} incurred many times more of them than the overloaded version.

\subsection{Allocator Design and Workload Fit}
\label{sec:rq-policy}

\subsubsection{Known Access Patterns and Heuristics}
\emph{When could a known access pattern make an allocator heuristic
effective?}
A heuristic is effective when the workload matches the assumption it encodes.
Knowing the pattern in advance lets you pick the heuristic that skips exactly
the work your workload would never need.

\subsubsection{System Allocator on Small Nodes}
\emph{Why might the system allocator outperform a custom allocator for small
nodes?}
Our linked-list test was the worst case for the custom allocator.
It was slower than regular \texttt{new} at every scale and had 2.55 times the
peak memory.
Our matrix-based tests differed by only 3--28\,\% in memory.
The cause is that per-allocation overhead is a fixed cost regardless of object
size, so it dominates when the objects are small.

\subsubsection{Fragmentation}
\emph{How does fragmentation affect allocation search time and memory use?}
Fragmentation means free memory exists but is split into pieces too small to
satisfy a request.
This means an allocation can fail even when total free space is
sufficient~\cite{wilson1995}.
Search time grows because there are more free-list entries to examine.

\subsubsection{Canaries and Allocation Metadata}
\emph{How do canaries and allocation metadata affect speed and memory use?}
In v7, \texttt{allocation\_overhead} is 64 bytes per allocation: a 48-byte
header, an 8-byte owner pointer, and two 4-byte canaries.
For the 72-byte linked-list node, header, alignment padding, and canaries bring
the block to 144 bytes, which the slab refill then rounds up to 192 bytes.
Every allocation writes two canaries and an owner pointer.
Every deallocation reads and compares both canaries before freeing.
This operation is cheap but can be significant when working with a small node.

\subsubsection{Cache Limits and Size Classes}
\emph{How should cache limits and size classes be selected for a workload?}
They should be chosen by measuring the workload's actual request sizes rather
than fixed in advance.
Cache limits should be sized to each thread's working set per class and capped
by a total byte budget multiplied across all threads.

\subsubsection{Workload-Specific Optimization}
\emph{When could an allocator optimized for one workload harm another
workload?}
Optimizations usually assume some specific scenario, so if that condition is
not met, the optimization can cause unnecessary overhead in another workload
compared to the ideal workload.

\subsection{Hardware and Operating System Effects}
\label{sec:rq-hardware}

\subsubsection{Matrix Calculation Time}
\emph{Why can matrix calculation time hide allocation improvements?}
Since the allocator versions do not change the matrix arithmetic, we look at
Amdahl's law~\cite{amdahl1967}.
It states that an improvement to allocation is capped by allocation's share of
total runtime.
Matrix multiplication has $O(n^3)$ arithmetic over $O(n^2)$ memory.
This means that as the matrix grows, the allocation share shrinks toward zero.

\subsubsection{Cache-Line Padding and False Sharing}
\emph{Several threads update adjacent objects. How could padding each object
to a cache-line boundary waste space yet improve performance by reducing false
sharing and cache coherence traffic?}
Cache coherence works on whole cache lines~\cite{hennessy2019}.
So, when two threads write to objects that share a line, each write
invalidates the other thread's copy and the line bounces between them.
Padding each object to a cache-line boundary wastes space, but it eliminates
that false sharing by trading cheap memory for expensive coherence traffic that
gets worse as thread count rises.

\subsection{Software Engineering Practice}
\label{sec:rq-engineering}

\subsubsection{Namespaces and Multiple Allocators}
\emph{How could namespaces support multiple allocators based on scope and
purpose?}
Namespaces let each subsystem use an allocator suited to its purpose.
Scoping allocators this way also isolates each subsystem's statistics and
failures, and a single namespace can swap a module's allocator without
changing its code.

\subsubsection{Custom Allocator Tradeoffs}
\emph{What tradeoffs come with creating and maintaining a custom allocator?}
A custom allocator can help maximize the efficiency of the program by reducing
the number of system calls and potentially memory fragmentation.
However, if the task is trivial, then the overhead from creating and
maintaining the custom allocator could slow down the entire process.
Overhead includes both runtime overhead and the overhead of writing, testing,
and maintaining the custom allocator.
The additional code can be complex and hard to read, which makes maintenance
harder than using system calls.

\subsubsection{Design, Portability, and Reusability}
\emph{Why do good design, separation of concerns, portability, and reusability
matter when code is shared across projects?}
Shared code must be readable and maintainable by people who did not write it.
This allows anyone to understand what a change will affect.
Portability was a constraint in this project because the various environments
handled the source code differently.

\subsubsection{Single Source of Truth}
\emph{Explain the importance of single source of truth in source code
management.}
A single source of truth means each piece of information is defined in exactly
one authoritative place.
This removes any ambiguity about which version is current when several people
edit the same file.
